%% kapa: exam=1 round=34 year=2023 period=2 form=A subject=appraisal_law subject_ko=감정평가관계법규 seq=16 number=16
%% source: Q-Net 감정평가사 기출문제 (cst00309) artlSeq=5211789 file=raw/1008/5211789/2023년도 제34회 감정평가사 1차시험 2교시.pdf
%% conversion: kordoc 4.24.51 -> tools/gen1.py (auto); review=proofread-sonnet; flags=none
\begin{question}{16}
부동산 가격공시에 관한 법령상 주택가격의 공시에 관한 설명으로 옳지 않은 것은?
\begin{choices}
  \item 이의신청의 기간·절차 및 방법은 표준주택가격을 공시할 때 관보에 공고해야 하는 사항이다.
  \item 표준주택가격의 공시에는 표준주택의 용도가 포함되어야 한다.
  \item 비주거용 집합부동산의 조사 및 산정 지침은 중앙부동산가격공시위원회의 심의 대상이다.
  \item 국토교통부장관은 표준주택가격을 조사·산정한 때에는 둘 이상의 감정평가법인등의 검증을 받아야 한다.
  \item 공동주택가격의 공시에는 공동주택의 면적이 포함되어야 한다.
\end{choices}
\end{question}
\answer{4}
